% ============================================================
% 3. Functional & Non-Functional Requirements
% ============================================================
\chapter{Requirements}\label{ch:requirements}

This chapter lists the functional and non-functional requirements of
CiboAmico. Every requirement is expressed with an active verb, is independent
from the graphical interface and is objectively testable.

\section{Functional Requirements}

\begin{reqs}
    \item \textbf{FR-01:} The system shall add a product to the household
    inventory with name, quantity, expiry date, storage position and unit of
    measure.
    \item \textbf{FR-02:} The system shall present the inventory ordered by
    ascending expiry date.
    \item \textbf{FR-03:} The system shall flag the products whose expiry date
    falls within the next three days.
    \item \textbf{FR-04:} The system shall provide the details of a product
    associated with a scanned barcode or a specific search query.
    \item \textbf{FR-05:} The system shall compute the list of ingredients
    missing from the inventory for a selected recipe.
    \item \textbf{FR-06:} The system shall publish a product on the
    marketplace only if the vendor is approved and the price is greater than
    zero.
    \item \textbf{FR-07:} The system shall create a single direct order for
    one buyer and one vendor, tracking the order state transitions.
    \item \textbf{FR-08:} The system shall create a recipe with at least two
    ingredients, published in the \emph{PROPOSED} state.
    \item \textbf{FR-09:} The system shall approve or suspend a vendor.
    \item \textbf{FR-10:} The system shall approve or reject a proposed
    recipe.
    \item \textbf{FR-11:} The system shall authenticate a user by e-mail and
    password, managing a session.
    \item \textbf{FR-12:} The system shall find the recipes compatible with the
    available inventory, including equivalent ingredient substitutions.
    \item \textbf{FR-13:} The system shall present to a vendor the orders
    received for their products, with their current state, and shall update
    the state upon the vendor's action.
\end{reqs}

\section{Business Rules}

Business rules constrain the functional behaviour:

\begin{reqs}
    \item \textbf{BR-01:} A product whose expiry date is in the past cannot be
    sold.
    \item \textbf{BR-02:} Only vendors in the \emph{APPROVED} state can
    publish products.
    \item \textbf{BR-03:} The available quantity of a product cannot be
    negative.
    \item \textbf{BR-04:} The order state machine restricts transitions
    exclusively to: CREATED $\rightarrow$ CONFIRMED/ANNULLED, CONFIRMED $
    \rightarrow$ IN\_DELIVERY/ANNULLED, IN\_DELIVERY $\rightarrow$ DELIVERED.
    \item \textbf{BR-05:} A recipe must contain at least two ingredients.
    \item \textbf{BR-06:} The price of a marketplace product must be greater
    than zero.
\end{reqs}

\section{Non-Functional Requirements}

\begin{reqs}
    \item \textbf{NFR-01:} The system shall support at least three persistence
    modes (relational database, local JSON files, in-memory demo), selectable
    at application startup.
    \item \textbf{NFR-02:} The system shall protect data access against
    SQL-injection by using prepared statements.
    \item \textbf{NFR-03:} The system shall store user passwords using a
    one-way cryptographic hash with salt.
    \item \textbf{NFR-04:} The line coverage of the automated test suite shall
    be at least 80\% on the application logic.
\end{reqs}

\section{Requirements--Use-Case Traceability}

Table~\ref{tab:trace} maps every functional requirement to the use cases that
realise it, guaranteeing full traceability from requirements to code.

\begin{table}[H]
\centering
\caption{FR to UC traceability.}
\label{tab:trace}
\begin{tabular}{@{}ll@{}ll@{}ll@{}ll@{}ll@{}}
\toprule
\textbf{FR} & \textbf{UC} & \textbf{FR} & \textbf{UC} & \textbf{FR} & \textbf{UC} \\
\midrule
FR-01 & UC-01 & FR-05 & UC-03 & FR-09 & UC-08 \\
FR-02 & UC-01 & FR-06 & UC-05 & FR-10 & UC-08 \\
FR-03 & UC-01 & FR-07 & UC-04 & FR-11 & UC-11 \\
FR-04 & UC-01 & FR-08 & UC-07 & FR-12 & UC-02 \\
 & & & & FR-13 & UC-06 \\
\bottomrule
\end{tabular}
\end{table}
